\documentclass[10pt,conference]{IEEEtran}

% ── Packages ─────────────────────────────────────────────────────────────────
\usepackage[utf8]{inputenc}
\usepackage[T1]{fontenc}
\usepackage{amsmath,amssymb,bm}
\usepackage{graphicx}
\usepackage{booktabs}
\usepackage{array}
\usepackage{multirow}
\usepackage{xcolor}
\usepackage{hyperref}
\usepackage{enumitem}
\usepackage{caption}
\usepackage{float}
\usepackage{cite}
\usepackage{url}

% ── Hyperref setup ────────────────────────────────────────────────────────────
\hypersetup{
  colorlinks=true,
  linkcolor=blue!60!black,
  citecolor=green!50!black,
  urlcolor=blue!70!black,
  pdftitle={Multimodal Video Summarization via Supervised Learning},
}

% ── Math operators ────────────────────────────────────────────────────────────
\DeclareMathOperator{\sigmoid}{\sigma}

% =============================================================================
\begin{document}

% ── Title ─────────────────────────────────────────────────────────────────────
\title{Multimodal Video Summarization via Supervised Learning:\\
Multimodal BiLSTM with Temporal Attention\\
{\large and Evaluation of the VideoIntel Content Analysis Application}}

\author{
  \IEEEauthorblockN{Tran Hai Duc}
  \IEEEauthorblockA{
    Student ID: 23127173\\
    University of Science, VNU-HCM\\
    thduc23@clc.fitus.edu.vn
  }
  \and
  \IEEEauthorblockN{Nguyen Cong Chien}
  \IEEEauthorblockA{
    Student ID: 23127331\\
    University of Science, VNU-HCM\\
    ncchien23@clc.fitus.edu.vn
  }
}

\maketitle

% ── Abstract ──────────────────────────────────────────────────────────────────
\begin{abstract}
This report presents a supervised video summarization method evaluated on the
SumMe and TVSum benchmarks. Visual features are extracted with a ResNet-50
backbone, while semantic speech features are obtained from Whisper automatic
speech recognition and SentenceBERT sentence encodings. A two-layer
Bidirectional LSTM (BiLSTM) combined with a Temporal Attention mechanism
produces per-frame importance scores. Inference follows a fixed summary ratio
of $r{=}0.35$ with Gaussian smoothing, shot segmentation, and boundary
expansion to assemble the final dynamic summary. A complementary evaluation
uses the VideoIntel application---built on CLIP, Whisper, and a local
LLM---to measure the retention of search and question-answering capabilities
when the original video is replaced by its summary. All content is
self-contained; no external source files are required.
\end{abstract}

\begin{IEEEkeywords}
video summarization, supervised learning, BiLSTM, temporal attention,
multimodal, SumMe, TVSum, CLIP, Whisper, video question answering
\end{IEEEkeywords}

% =============================================================================
\section{Introduction}
% =============================================================================

\subsection{The Video Summarization Problem}

\textbf{Video summarization} aims to produce a condensed representation of an
original video that preserves its key content and narrative flow. Two common
output forms are distinguished:
\begin{itemize}
  \item \textbf{Static summaries:} a selected set of representative keyframes
        or key shots.
  \item \textbf{Dynamic summaries:} a shorter video file assembled from
        selected segments, with audio preserved in sync.
\end{itemize}

The user-facing goal is to \emph{reduce viewing time} while retaining most of
the important information; the system-level goal is to \emph{reduce data
dimensionality in a structured manner} rather than compressing uniformly over
time.

\subsection{Methodological Context}

Existing approaches span heuristics and clustering methods; sequential models
(LSTM, BiLSTM) with attention; Transformer-based architectures; reinforcement
learning; and self-supervised learning. \textbf{Multimodal} methods that
jointly process vision and audio/speech signals are particularly effective for
semantically rich segments with little visual change---for example, lecture
slides accompanied by narration.

\subsection{Contributions}

The specific contributions of this report are:
\begin{enumerate}
  \item A \textbf{multimodal supervised video summarization} pipeline trained
        on \textbf{SumMe} and \textbf{TVSum}, combining visual and speech
        features.
  \item A \textbf{BiLSTM} architecture with \textbf{Temporal Attention}
        operating over 2432-dimensional multimodal feature sequences; the
        speech branch uses \textbf{Whisper} (ASR) and \textbf{SentenceBERT}.
  \item A \textbf{dynamic video summary} $V_{\mathrm{sum}}$ with synchronized
        audio assembled via \texttt{ffmpeg} at a fixed summary ratio
        $r = 0.35$.
  \item An auxiliary evaluation through the \textbf{VideoIntel} application,
        which measures \textbf{semantic search} (CLIP) and
        \textbf{question-answering} (Whisper + local LLM) quality on
        original/summary video pairs.
\end{enumerate}

\subsection{Related Work}

General surveys~\cite{haq2021,money2015,apostolidis2021survey,hu2023,kadam2022}
provide a broad overview of classical and deep-learning pipelines. The standard
pipeline comprises \textbf{frame-level feature extraction}, \textbf{importance
score estimation}, and \textbf{segment selection} under a time budget.

For supervised deep-learning methods, \textbf{LSTM/BiLSTM} variants such as
vsLSTM and dppLSTM~\cite{zhang2016lstm} model temporal dependencies;
\textbf{attention} mechanisms (VASNet~\cite{fajtl2018}, Temporal Attention in
this work) highlight informationally dense temporal regions. The foundational
\textbf{Transformer} architecture~\cite{vaswani2017} based on self-attention
has inspired encoder-only and full encoder-decoder variants for video
summarization~\cite{lan2025fulltransnet,ji2019avs}. Unsupervised and
adversarial directions (SUM-GAN~\cite{mahasseni2017},
SUM-GAN-AED~\cite{apostolidis2023sumganaed}, Cycle-SUM~\cite{zhang2019})
and Transformer-based models (MSVA~\cite{ghauri2021}, PGL-SUM~\cite{apostolidis2021pgl},
MHSCNet~\cite{zhu2022}) provide additional architectural alternatives.
\textbf{Graph Neural Networks} applied to spatial-temporal
reasoning~\cite{zhu2022rrstg} capture object-level dependencies that
frame-level methods miss. \textbf{Multimodal} fusion of vision and speech
is an active research area~\cite{psallidas2021,he2023align}; language-guided
summarization via CLIP-based attention~\cite{narasimhan2021clipit} shows
strong zero-shot generalization.

At the feature level, \textbf{CNN ResNet}~\cite{he2016} provides visual
vectors; the text branch uses \textbf{Whisper}~\cite{radford2023whisper} and
\textbf{SentenceBERT}~\cite{reimers2019} for fixed-dimension fusion.
\textbf{CLIP}~\cite{radford2021clip} provides a joint image-text embedding
space for semantic search. Video recognition architectures
(TimeSformer~\cite{bertasius2021}, Video Swin~\cite{liu2022swin},
I3D~\cite{carreira2017}) are relevant as downstream evaluation references.
Comprehensive surveys cover deep learning methods for video
summarization~\cite{apostolidis2021survey,yu2026deeplearning,tiwari2021survey,meena2023review}.

% =============================================================================
\section{Problem Formulation}
% =============================================================================

\subsection{Input}
\begin{itemize}
  \item An input video $V$ in a standard format, of arbitrary duration.
  \item During training: for each sampled time step $t \in \{1,\dots,T\}$, a
        scalar importance label $y_t \in [0,1]$ derived from SumMe/TVSum
        annotations (normalized according to the dataset convention).
\end{itemize}

\subsection{Output}
\begin{itemize}
  \item \textbf{Static summary (auxiliary):} a set of selected keyframes or
        shots.
  \item \textbf{Dynamic summary (primary):} a video $V_{\mathrm{sum}}$ with
        total duration approximately $r \cdot |V|$, where $r = 0.35$ is fixed
        at inference time. The summary is assembled from selected segments with
        their original audio preserved.
\end{itemize}

\subsection{Deployment Constraints}

The \textbf{dynamic summary} is the primary output, as it best matches
real-world user needs (fast viewing with audio). Audio extraction and
\texttt{ffmpeg}-based assembly are used throughout the pipeline.

% =============================================================================
\section{Theoretical Background and Mathematical Formulation}
% =============================================================================

\subsection{Multimodal Feature Space}

At each sampled time step $t$ (e.g., at 2\,fps), we construct a feature
vector:
\begin{equation}
  \mathbf{x}_t =
  \begin{bmatrix}
    \mathbf{v}_t \\ \mathbf{a}_t
  \end{bmatrix}
  \in \mathbb{R}^{D_v + D_a},
  \quad D_v = 2048,\; D_a = 384,\; D_v + D_a = 2432.
\end{equation}
Here $\mathbf{v}_t$ is the CNN visual feature (ResNet-50) and $\mathbf{a}_t$
is the temporally aligned semantic speech feature (Section~\ref{sec:audio}).

\subsection{Supervised Learning and Binary Labels}

The model produces a per-frame logit $z_t \in \mathbb{R}$. The estimated
probability is $\hat{y}_t = \sigma(z_t)$, where $\sigma$ is the sigmoid
function. Ground-truth labels $y_t \in [0,1]$ are \textbf{binarized} at
threshold~0.5 to obtain $\tilde{y}_t \in \{0,1\}$ for loss computation
(BCEWithLogitsLoss). The \textbf{loss function} averaged over $T$ steps is:
\begin{equation}
  \mathcal{L}
  = -\frac{1}{T}\sum_{t=1}^{T}
  \Bigl[
    \tilde{y}_t \log \sigma(z_t)
    + (1 - \tilde{y}_t)\log\bigl(1 - \sigma(z_t)\bigr)
  \Bigr].
  \label{eq:bce}
\end{equation}

\subsection{Bidirectional LSTM (BiLSTM)}

Given input sequence $\mathbf{X} = (\mathbf{x}_1,\dots,\mathbf{x}_T)$, a
forward LSTM produces $\overrightarrow{\mathbf{h}}_t \in \mathbb{R}^{H}$ and a
backward LSTM produces $\overleftarrow{\mathbf{h}}_t \in \mathbb{R}^{H}$,
with $H = 256$. The concatenated hidden state is:
\begin{equation}
  \mathbf{h}_t = \bigl[\overrightarrow{\mathbf{h}}_t \;;\;
                       \overleftarrow{\mathbf{h}}_t\bigr]
               \in \mathbb{R}^{2H} = \mathbb{R}^{512}.
\end{equation}
Two stacked BiLSTM layers deepen the representation; each layer reads the
sequence in \textbf{both temporal directions}, so the importance estimate at
time $t$ depends on both past and future context.

\subsection{Temporal Attention}

Let $\mathbf{H} = [\mathbf{h}_1,\dots,\mathbf{h}_T] \in \mathbb{R}^{T \times 512}$.
Attention scores and weights are computed as:
\begin{equation}
  e_t = \mathbf{u}^\top \tanh(\mathbf{W}_a \mathbf{h}_t + \mathbf{b}_a),
  \quad
  \alpha_t = \frac{\exp(e_t)}{\sum_{k=1}^{T}\exp(e_k)}.
\end{equation}
A globally pooled context vector is:
\begin{equation}
  \mathbf{c} = \sum_{t=1}^{T} \alpha_t\, \mathbf{h}_t.
\end{equation}
To retain a \textbf{per-frame logit} for shot selection, $\mathbf{h}_t$ is
combined with $\mathbf{c}$ via concatenation or a linear projection.

\subsection{Output Logit Layer}

\begin{equation}
  z_t = \mathbf{w}_o^\top \mathbf{h}_t^{\mathrm{out}} + b_o,
\end{equation}
where $\mathbf{h}_t^{\mathrm{out}}$ is the post-attention representation.
Sigmoid is applied at inference as $\hat{y}_t = \sigma(z_t)$.

\subsection{Datasets and Evaluation Metrics}

\begin{itemize}
  \item \textbf{SumMe}~\cite{gygli2014}: per-frame importance labels from
        multiple annotators.
  \item \textbf{TVSum}~\cite{song2015}: per-segment importance labels from
        multiple annotators.
  \item \textbf{F-score}: harmonic mean of precision and recall over selected
        vs.\ ground-truth frames/segments.
  \item \textbf{Precision / Recall}: fraction of selected frames that are
        truly important / fraction of important frames that are selected.
  \item \textbf{Temporal overlap}: temporal intersection ratio between
        predicted and ground-truth segments.
\end{itemize}

% =============================================================================
\section{Pipeline: Model Architecture and Forward Pass}
\label{sec:pipeline}
% =============================================================================

\subsection{Module Overview}

The data-flow pipeline is as follows:
\begin{enumerate}
  \item \textbf{Frame sampling} at 2\,fps; sequence length capped at
        \texttt{max\_seq\_len}$\,{=}\,960$ frames ($\approx$8\,min).
  \item \textbf{Visual backbone:} ResNet-50 extracts
        $\mathbf{v}_t \in \mathbb{R}^{2048}$.
  \item \textbf{Audio-language branch:} \texttt{ffmpeg} $\to$ Whisper $\to$
        SentenceBERT $\to$ $\mathbf{a}_t$ (Section~\ref{sec:audio}).
  \item \textbf{Feature concatenation:}
        $\mathbf{x}_t = [\mathbf{v}_t;\, \mathbf{a}_t] \in \mathbb{R}^{2432}$.
  \item \textbf{Two-layer BiLSTM}, hidden size 256 per direction
        $\to$ $\mathbf{h}_t \in \mathbb{R}^{512}$.
  \item \textbf{Temporal Attention} $\to$ linear head $\to$ $z_t$.
  \item Loss: \textbf{BCEWithLogitsLoss}; optimizer: \textbf{AdamW};
        scheduler: \textbf{StepLR}.
\end{enumerate}

\subsection{Tensor Shapes and Batching}

Sequences longer than 960 steps are \textbf{truncated}. Batches use
\textbf{padding and masking} so that loss is not computed over padded
positions. Each valid time step carries a 2432-dimensional input.

\subsection{Ablation Variants}

\begin{table}[h]
\centering
\caption{Ablation Variants and Architectural Differences}
\label{tab:ablations}
\begin{tabular}{ll}
\toprule
\textbf{Variant} & \textbf{Modification} \\
\midrule
Full model   & $\mathbf{x}_t = [\mathbf{v}_t;\mathbf{a}_t]$, with Temporal Attention \\
Visual-only  & Drop $\mathbf{a}_t$; reduced input dimension \\
No Attention & Remove Temporal Attention; BiLSTM + linear head only \\
UniLSTM      & \texttt{bidirectional=False}; single-direction LSTM \\
\bottomrule
\end{tabular}
\end{table}

% =============================================================================
\section{Pipeline: Audio Branch, Text Encoding, and Temporal Alignment}
\label{sec:audio}
% =============================================================================

\subsection{Audio Extraction}

\texttt{ffmpeg} converts the video's audio track to \textbf{WAV} format
(mono, 16\,kHz), synchronized to the video timeline.

\subsection{Automatic Speech Recognition (Whisper)}

Whisper~\cite{radford2023whisper} produces \textbf{sentence/segment tokens}
each carrying timestamps $(t_k^{\mathrm{start}}, t_k^{\mathrm{end}})$ and raw
transcript text.

\subsection{SentenceBERT Encoding}

Each transcript segment is encoded into
$\mathbf{s}_k \in \mathbb{R}^{384}$ using SentenceBERT~\cite{reimers2019}.

\subsection{Frame-Level Alignment}

For each sampled frame at absolute time $T_t$:
\begin{itemize}
  \item If $T_t \in [t_k^{\mathrm{start}}, t_k^{\mathrm{end}}]$, assign
        $\mathbf{a}_t = \mathbf{s}_k$.
  \item Otherwise: assign $\mathbf{a}_t = \mathbf{0}$ or interpolate from
        the nearest segment.
\end{itemize}

\subsection{Feature Fusion}

\begin{equation}
  \mathbf{x}_t = [\mathbf{v}_t;\, \mathbf{a}_t] \in \mathbb{R}^{2432}.
\end{equation}
The audio-language branch increases the weight of \textbf{semantically rich
speech segments} even when the visual signal is relatively static.

% =============================================================================
\section{Pipeline: Training, Inference, and Video Assembly}
% =============================================================================

\subsection{Training (Offline)}

\begin{enumerate}
  \item \textbf{Pre-processing:} sample at 2\,fps; ResNet-50 $\to$ visual
        tensors; Whisper + SentenceBERT $\to$ aligned audio-text features.
  \item \textbf{Construct} $\mathbf{x}_t$ sequences; build padded/masked
        batches.
  \item \textbf{Forward pass:} $\mathbf{x}_{1:T} \to$ BiLSTM $\to$ Attention
        $\to$ $z_{1:T}$.
  \item \textbf{Backward pass:} compute $\mathcal{L}$ (Eq.~\ref{eq:bce});
        update with AdamW; apply StepLR; early stopping; save
        \texttt{best.pt}.
\end{enumerate}

\subsection{Inference on a New Video}

\begin{enumerate}
  \item Extract $\{\mathbf{v}_t\}$ and $\{\mathbf{a}_t\}$ as in training.
  \item Forward pass $\to$ $\hat{y}_t = \sigma(z_t)$.
  \item \textbf{Gaussian smoothing} to reduce temporal noise.
  \item \textbf{Shot segmentation} via \texttt{keyshot.py}.
  \item \textbf{Shot selection} such that total duration $\approx r \cdot
        |V|$, $r = 0.35$.
  \item \textbf{Boundary expansion} of $\pm 1.5$\,s per selected segment;
        merge nearby segments if needed.
  \item \texttt{ffmpeg} cuts and concatenates segments into
        $V_{\mathrm{sum}}$ with full audio.
\end{enumerate}

% =============================================================================
\section{Hyperparameters and Experimental Configuration}
% =============================================================================

\begin{table}[h]
\centering
\caption{Hyperparameters and Experimental Configuration}
\label{tab:hyperparams}
\begin{tabular}{ll}
\toprule
\textbf{Component} & \textbf{Value / Notes} \\
\midrule
Architecture         & Multimodal BiLSTM + Temporal Attention \\
Visual backbone      & ResNet-50, $D_v = 2048$ \\
Audio-text features  & Whisper + SentenceBERT, $D_a = 384$ \\
Input dimension      & $2432 = 2048 + 384$ \\
Frame sampling       & 2\,fps, \texttt{max\_seq\_len} = 960 \\
Smoothing            & Gaussian, $\sigma$ from config file \\
Summary ratio        & $r = 0.35$ \\
Boundary expansion   & $\pm 1.5$\,s \\
Loss                 & BCEWithLogitsLoss \\
Optimizer            & AdamW, lr $= 10^{-3}$, wd $= 10^{-5}$ \\
Scheduler            & StepLR (step$\,{=}\,15$, $\gamma{=}0.5$) \\
\bottomrule
\end{tabular}
\end{table}

% =============================================================================
\section{Experimental Results and Ablation Study}
\label{sec:ablation}
% =============================================================================

\subsection{Ablation Study Design}

An ablation study evaluates a model by \emph{holding the dataset, train/test
split, and evaluation protocol fixed} while \emph{enabling or disabling
individual architectural blocks}---for instance, dropping a feature branch,
removing attention, or switching the LSTM directionality. Comparing metrics
before and after each removal answers the question: \emph{does this component
carry useful signal, or does it merely add complexity?}

Four variants are examined, motivated as follows.

\textbf{Full model (Multimodal BiLSTM + Temporal Attention).}
This is the complete proposed architecture: visual features from ResNet-50 are
concatenated with SentenceBERT-encoded Whisper transcript embeddings, and a
two-layer BiLSTM processes the joint sequence under a Temporal Attention
mechanism.

\textbf{Visual-only (no audio branch).}
Video summarization often relies primarily on visual content; however, for
speech-rich material such as TVSum news clips, the language branch may supply
contextual signal that is not recoverable from frames alone. Dropping
$\mathbf{a}_t$ isolates the \emph{marginal contribution of audio/text} relative
to a vision-only CNN baseline~\cite{zhang2016lstm}.

\textbf{No Temporal Attention.}
The attention module aggregates the global context vector $\mathbf{c}$, which
enables the frame-level scorer to compare each position against the full
sequence distribution. Removing attention reduces the architecture to a
BiLSTM followed by a direct linear scorer, testing the
\emph{necessity of sequence-level context pooling}~\cite{fajtl2018}.

\textbf{UniLSTM (unidirectional).}
A BiLSTM provides each position with both past and future context, which is
particularly valuable for the importance-scoring task where a shot's relevance
can only be judged relative to what comes later. Replacing
\texttt{bidirectional=True} with \texttt{bidirectional=False} isolates the
\emph{gain from bidirectional encoding}~\cite{zhang2016lstm,apostolidis2021pgl}.

All four variants are trained with identical hyperparameters (Table~\ref{tab:hyperparams}),
the same random seed, and the same post-processing pipeline ($r{=}0.35$,
Gaussian smoothing, boundary expansion $\pm 1.5$\,s).

\subsection{Evaluation Protocol}

Ground-truth labels are binarized using the \emph{same top-$k$ ratio}
as the predicted summaries (controlled by \texttt{summary\_ratio}). As a
direct consequence, the number of positive predictions and the number of
positive ground-truth frames are approximately equal, which causes
\textbf{Precision and Recall to converge toward the same value} in the
aggregate---this is a property of the evaluation protocol, not a numerical
coincidence. F-score is therefore the harmonic mean of two nearly identical
values and reduces to approximately Precision~$\approx$~Recall for these
runs. Temporal overlap is reported separately as a frame-level IoU between
predicted and ground-truth segments, and is \emph{not} affected by the
top-$k$ binarization symmetry.

SumMe and TVSum scores are kept separate throughout, as the two datasets
differ substantially in video length, annotation style, and domain; merging
them into a single number would obscure dataset-specific behavior.

\subsection{Quantitative Ablation Results}

\begin{table}[h]
\centering
\caption{Ablation Study: F-Score, Precision/Recall, and Temporal Overlap (\%)}
\label{tab:ablation_results}
\setlength{\tabcolsep}{4pt}
\begin{tabular}{p{3.0cm}cccc}
\toprule
\textbf{Variant} & \textbf{SumMe} & \textbf{TVSum} & \textbf{P/R} & \textbf{Ovlp} \\
 & \textbf{F (\%)} & \textbf{F (\%)} & \textbf{(\%)} & \textbf{(\%)} \\
\midrule
\textbf{Full model (Ours)}  & 26.9 & 60.1 & 47.7 & 23.5 \\
Visual-only                 & 31.4 & 59.7 & 49.1 & 33.3 \\
No Attn.                    & 34.9 & 63.7 & 52.9 & 27.4 \\
UniLSTM                     & 35.4 & 63.1 & 52.7 & 37.3 \\
\bottomrule
\end{tabular}
\end{table}

\subsection{Analysis and Interpretation}

Several observations merit discussion.

\textbf{Dataset asymmetry.}
TVSum F-scores consistently exceed SumMe scores across all variants by
a substantial margin ($\approx$27--28\,pp). This reflects
fundamental differences between the two benchmarks: TVSum videos are
longer and their importance labels---derived from web-video titles---tend
to be more concentrated, making top-$k$ selection easier. SumMe uses
crowdsourced per-frame scores with higher annotator disagreement, which is
known to increase evaluation variance~\cite{gygli2014,song2015}. Results on
one benchmark should not be used to predict absolute performance on the other.

\textbf{Full model does not rank first on SumMe.}
The full multimodal model achieves lower SumMe F-score than the
visual-only, no-attention, and UniLSTM variants in this run. This outcome
is consistent with known phenomena in multimodal learning: a richer
architecture has more parameters and requires more careful regularization,
learning-rate tuning, and early stopping; with a fixed training budget the
simpler variants may converge to a better local optimum on a smaller
dataset such as SumMe ($<25$ videos in the standard split)~\cite{apostolidis2021pgl}.
Additionally, if the audio transcripts in the SumMe subset are sparse or
misaligned (SumMe contains user-generated short videos, many without
sustained narration), the audio branch introduces noise rather than
signal. A rigorous comparison would require multiple seeds with statistical
significance testing; the present single-seed result should be interpreted
as a \emph{specific reproducible outcome} rather than a universal law.

\textbf{Precision equals Recall in the aggregate Prec./Rec.\ column.}
As explained in Section~\ref{sec:ablation}.B, the top-$k$ binarization
symmetry of the evaluator causes the aggregate positive counts of
predictions and ground-truth to be nearly identical, so Precision $\approx$
Recall $\approx$ F-score in the combined metric column. The per-dataset
SumMe and TVSum F-scores are computed from the respective subsets and are
not subject to this equivalence.

\textbf{Temporal overlap.}
The UniLSTM variant achieves the highest temporal overlap (37.3\%),
followed by visual-only (33.3\%). The full model achieves the lowest
overlap (23.5\%). One interpretation is that the full model's attention
mechanism concentrates high-score predictions in fewer, shorter temporal
regions, yielding a sparser but more discriminative summary that overlaps
less with the broad ground-truth intervals. An alternative explanation is
that the audio branch shifts score peaks toward speech-dense regions that
do not coincide with the annotated shot boundaries in this split.
Both effects could be disentangled by examining score histograms and
shot-length distributions, which are left for future work.

\textbf{Implications for architecture design.}
The ablation results indicate that \emph{architecture design choices
interact with dataset characteristics} in non-trivial ways. The audio
branch improves TVSum performance (60.1 vs.\ 59.7 for visual-only,
a difference of +0.4\,pp on a semantically richer set), while
apparently hurting SumMe ($-4.5$\,pp). Temporal Attention and
bidirectionality follow a similar pattern. These findings suggest that
for \emph{mixed-content or narration-heavy} corpora, multimodal
features are beneficial, whereas for \emph{visually dynamic, low-speech}
content, simpler unimodal models may be preferable. A dataset-adaptive
weighting of modalities---for example, scaling $\mathbf{a}_t$ by an
estimated speech-activity ratio---is a natural extension.

\subsection{Comparison with Published Methods}

\begin{table}[h]
\centering
\caption{Comparison with Published Methods on SumMe and TVSum (F-Score, \%)}
\label{tab:literature}
\begin{tabular}{lllcc}
\toprule
\textbf{Method} & \textbf{Year} & \textbf{Type} & \textbf{SumMe} & \textbf{TVSum} \\
\midrule
vsLSTM~\cite{zhang2016lstm}       & 2016 & Supervised   & 37.6 & 54.2 \\
dppLSTM~\cite{zhang2016lstm}      & 2016 & Supervised   & 38.6 & 54.7 \\
SUM-GAN~\cite{mahasseni2017}      & 2017 & Unsupervised & 38.7 & 50.8 \\
Cycle-SUM~\cite{zhang2019}        & 2019 & Unsupervised & 41.9 & 57.6 \\
VASNet~\cite{fajtl2018}           & 2018 & Supervised   & 49.6 & 61.4 \\
DSNet~\cite{zhu2021dsnet}         & 2020 & Supervised   & 51.2 & 61.9 \\
MSVA~\cite{ghauri2021}            & 2021 & Supervised   & 53.4 & 61.5 \\
PGL-SUM~\cite{apostolidis2021pgl} & 2021 & Supervised   & 57.1 & 62.7 \\
SUM-GDA                           & 2020 & ---          & 56.3 & 60.7 \\
\midrule
\textbf{This work (full model)}   & ---  & Supervised   & \textbf{26.9} & \textbf{60.1} \\
\bottomrule
\end{tabular}
\end{table}

Direct comparison should be interpreted cautiously due to differences in
dataset splits, summary ratios, and post-processing strategies.
The TVSum F-score of 60.1\% is competitive with vsLSTM (54.2\%) and
approaches VASNet (61.4\%), while SumMe performance remains below
published baselines, consistent with the dataset-sensitivity discussion
above.

\subsection{Comparison with Initial Target Metrics}

The initial targets were $\ge 0.43$ F-score on SumMe and $\ge 0.55$ on
TVSum. The full-model TVSum result of \textbf{0.601} surpasses the TVSum
threshold. The SumMe result of 0.269 falls below 0.43, reflecting the
strong sensitivity of SumMe performance to dataset split, training budget,
and audio-branch noise on short user-generated videos.

% =============================================================================
\section{VideoIntel Content Analysis Application}
% =============================================================================

\subsection{Motivation and Operating Hypotheses}

The summarization pipeline produces a \textbf{shorter, lighter} file, but may:
\begin{enumerate}
  \item \textbf{Reduce cost:} transcription time decreases with shorter audio.
  \item \textbf{Reduce information:} loss of speech segments or off-topic
        selection $\to$ \textbf{semantic drift}.
\end{enumerate}
VideoIntel measures both effects on \textbf{(original, summary) video pairs}
using: indexing $\to$ Search $\to$ Q\&A.

Technical assumptions: \textbf{CLIP} (ViT-B/32) provides a shared embedding
space for text and frames; \textbf{Whisper} provides transcripts; the
\textbf{local LLM (Ollama)} relies solely on the transcript---a short or
off-topic transcript can cause \textbf{hallucination}.

\subsection{Search Task}

The scoring scheme is:
\begin{equation}
  \mathbf{q} = \mathrm{CLIP}_{\text{text}}(q),\quad
  \mathbf{v}_i = \mathrm{CLIP}_{\text{image}}(\text{frame}_i),\quad
  \mathrm{score}_i = \mathbf{q} \cdot \mathbf{v}_i.
\end{equation}
Three fixed English queries benchmark CLIP top-1 score to estimate
\textbf{visual representativeness}: ``Person Speaking'' (PS), ``Important
Event'' (IE), and ``Diagram or Chart'' (DC).

\subsection{Q\&A Task}

A fixed question is posed to every video: \textit{``Please describe the
overall content of this video based on the transcript. Answer in Vietnamese,
approximately 3--5 sentences.''} Generation is capped at 512 tokens for fair
comparison.

\subsection{Benchmark Metric Definitions}

\begin{itemize}
  \item $\mathrm{duration\_ratio} = T_{\mathrm{sum}} / T_{\mathrm{orig}}$
  \item File size reduction: $1 - \mathrm{size}_{\mathrm{sum}} /
        \mathrm{size}_{\mathrm{orig}}$
  \item Transcript coverage: $\sum \Delta t$ of Whisper segments as fraction
        of video duration
  \item Segment retention ratio:
        $N_{\mathrm{seg,summ}} / N_{\mathrm{seg,orig}}$
  \item CLIP top-1: cosine similarity between query and best-matching frame
  \item Whisper speedup: $t_{\mathrm{orig}} / t_{\mathrm{sum}}$
\end{itemize}

\subsection{Quantitative Results (10 Video Pairs)}

\begin{table}[H]
\centering
\caption{Average VideoIntel Benchmark Results (10 Video Pairs)}
\label{tab:vi_avg}
\begin{tabular}{lc}
\toprule
\textbf{Metric} & \textbf{Value} \\
\midrule
Duration ratio (summary / original)             & $\approx$21.5\%  \\
File size reduction                              & $\approx$76.8\%  \\
Whisper speedup                                  & $\approx$4.4$\times$ \\
Transcript coverage (original / summary)         & $\approx$93.8\% / 94.2\% \\
Segment retention ratio                          & $\approx$25.6\%  \\
CLIP top-1 difference (summ $-$ orig)            & +0.0015 (negligible) \\
Acceptable Q\&A (qualitative)                    & Original 8/10, Summary 5/10 \\
\bottomrule
\end{tabular}
\end{table}

\begin{table}[H]
\centering
\caption{Duration and File Size Per Video Pair}
\label{tab:vi_size}
\begin{tabular}{lrrrrrr}
\toprule
\textbf{Video} & \textbf{Orig (s)} & \textbf{Summ (s)} & \textbf{Ratio} &
\textbf{Orig (MB)} & \textbf{Summ (MB)} & \textbf{Reduc.} \\
\midrule
video\_01 & 351  & 76  & 21.7\% & 17.22 & 3.80  & 77.9\% \\
video\_02 & 1198 & 254 & 21.2\% & 75.29 & 15.90 & 78.9\% \\
video\_03 & 480  & 104 & 21.7\% & 25.86 & 4.75  & 81.6\% \\
video\_04 & 488  & 111 & 22.7\% & 21.48 & 8.76  & 59.2\% \\
video\_05 & 1555 & 342 & 22.0\% & 73.16 & 16.06 & 78.0\% \\
video\_06 & 927  & 193 & 20.8\% & 48.01 & 11.15 & 76.8\% \\
video\_07 & 1359 & 314 & 23.1\% & 71.41 & 16.05 & 77.5\% \\
video\_08 & 831  & 177 & 21.3\% & 49.93 & 10.43 & 79.1\% \\
video\_09 & 1054 & 210 & 19.9\% & 60.34 & 11.83 & 80.4\% \\
video\_10 & 1006 & 209 & 20.8\% & 54.61 & 11.77 & 78.4\% \\
\midrule
\textbf{Avg.} & 825 & 179 & 21.5\% & 50.6 & 11.1 & 76.8\% \\
\bottomrule
\end{tabular}
\end{table}

\begin{table}[H]
\centering
\caption{Whisper Transcription and CLIP Encoding Times (60 Keyframes)}
\label{tab:vi_time}
\begin{tabular}{lrrrrr}
\toprule
\textbf{Video} & \textbf{W-Orig (s)} & \textbf{W-Summ (s)} &
\textbf{Speedup} & \textbf{C-Orig (s)} & \textbf{C-Summ (s)} \\
\midrule
video\_01 & 109.29 & 50.52 & $2.2\times$ & 3.03 & 3.15 \\
video\_02 & 112.15 & 25.20 & $4.5\times$ & 3.45 & 3.32 \\
video\_03 &  97.64 & 15.91 & $6.1\times$ & 3.46 & 3.44 \\
video\_04 & 207.15 & 29.73 & $7.0\times$ & 3.16 & 3.09 \\
video\_05 & 132.18 & 36.23 & $3.6\times$ & 3.20 & 3.12 \\
video\_06 & 155.04 & 44.66 & $3.5\times$ & 3.31 & 3.21 \\
video\_07 & 194.76 & 44.11 & $4.4\times$ & 3.12 & 3.32 \\
video\_08 & 117.31 & 25.34 & $4.6\times$ & 3.07 & 2.99 \\
video\_09 & 148.50 & 33.30 & $4.5\times$ & 3.26 & 3.34 \\
video\_10 & 155.26 & 38.14 & $4.1\times$ & 3.13 & 3.13 \\
\midrule
\textbf{Avg.} & 142.9 & 34.3 & $4.4\times$ & 3.22 & 3.21 \\
\bottomrule
\end{tabular}
\end{table}

\begin{table}[H]
\centering
\caption{Transcript Coverage and Segment Retention}
\label{tab:vi_seg}
\begin{tabular}{lccrrc}
\toprule
\textbf{Video} & \textbf{Cov Orig} & \textbf{Cov Summ} &
\textbf{Seg Orig} & \textbf{Seg Summ} & \textbf{\% Ret.} \\
\midrule
video\_01 & 98.4\% & 100.0\% & 103 &  31 & 30.1\% \\
video\_02 & 93.7\% &  99.4\% & 478 &  71 & 14.9\% \\
video\_03 & 93.1\% &  93.0\% &  89 &  19 & 21.3\% \\
video\_04 & 91.5\% &  97.7\% &  84 &  42 & 50.0\% \\
video\_05 & 96.1\% &  93.5\% & 320 & 138 & 43.1\% \\
video\_06 & 94.1\% &  92.3\% & 202 &  51 & 25.2\% \\
video\_07 & 88.7\% &  94.0\% & 400 &  81 & 20.3\% \\
video\_08 & 93.4\% &  85.3\% & 164 &  33 & 20.1\% \\
video\_09 & 92.5\% &  89.7\% & 378 &  36 &  9.5\% \\
video\_10 & 96.8\% &  97.2\% & 214 &  46 & 21.5\% \\
\midrule
\textbf{Avg.} & 93.8\% & 94.2\% & 243 & 55 & 25.6\% \\
\bottomrule
\end{tabular}
\end{table}

\begin{table}[H]
\centering
\caption{Average CLIP Top-1 Scores: Original vs.\ Summary (Three Queries)}
\label{tab:clip_scores}
\begin{tabular}{lrrr}
\toprule
\textbf{Query} & \textbf{Original} & \textbf{Summary} & \textbf{Diff.} \\
\midrule
Person Speaking (PS)        & 0.2509 & 0.2524 & $+0.0015$ \\
Important Event (IE)        & 0.2354 & 0.2374 & $+0.0020$ \\
Diagram or Chart (DC)       & 0.2274 & 0.2271 & $-0.0003$ \\
\midrule
\textbf{Mean}               & 0.2379 & 0.2390 & $+0.0015$ \\
\bottomrule
\end{tabular}
\end{table}

The overall mean CLIP top-1 difference (summary $-$ original) is
\textbf{+0.0015}, which is \textbf{statistically negligible}.

\subsection{Qualitative Q\&A Results}

Ratings: \textbf{Accurate} (topic matches, no factually incorrect details),
\textbf{Partial} (correct direction but missing details or hallucinated extras),
\textbf{Incorrect} (off-topic or wrong content).

\textbf{Overall:} On the original transcript, \textbf{8/10} responses were
acceptable. On the summary transcript, only \textbf{5/10} were Accurate; the
remaining 5 fell into Partial or Incorrect.

\begin{table}[H]
\centering
\caption{Per-Video Qualitative Q\&A Rating (Summary Transcript)}
\label{tab:qa_detail}
\begin{tabular}{p{1.1cm}p{1.3cm}p{5.0cm}}
\toprule
\textbf{Video} & \textbf{Rating} & \textbf{Observation} \\
\midrule
video\_01 & Incorrect
  & Original: US--Israel--Iran news. Summary: completely different topic
    (antibiotics / comedy) --- severe content-selection error. \\
video\_02 & Accurate
  & Both cover Israel/Iran, energy, US social issues --- topic matches. \\
video\_03 & Partial
  & Correct sports direction (FA Cup) but missing specific team/player
    details. \\
video\_04 & Accurate
  & Both cover AI scenario planning and control risks --- well aligned. \\
video\_05 & Accurate
  & BCIs, biomaterials in summary matches original neural electronics
    content. \\
video\_06 & Partial
  & Correct geopolitical direction but adds ``Vietnamese traffic accident''
    absent from original (hallucination). \\
video\_07 & Partial
  & Correct on economy/China but adds ``Vietnam economic competition''
    which is unrelated. \\
video\_08 & Partial
  & Correct on Iran/Trump/China but adds off-topic Chinese real estate
    and New York protests. \\
video\_09 & Partial
  & International part approximately correct; domestic content missing
    (only $\approx$9.5\% of segments retained). \\
video\_10 & Accurate
  & Both cover Middle East maritime issues --- key topics preserved. \\
\bottomrule
\end{tabular}
\end{table}

\textbf{Interpretation.} Q\&A quality on the summary does \emph{not} track
CLIP search quality: transcript density and topic fidelity are the critical
factors. The two extremes are \textbf{video\_01} (complete topic drift) and
\textbf{video\_09} (9.5\% segment retention causing severe information loss).

\subsection{Limitations and Recommendations}

\begin{itemize}
  \item Sample size: 10 videos; one LLM; one question template.
  \item CLIP top-1 score does not substitute for human-labeled relevance.
  \item \textbf{Recommendations:} (i) enforce a minimum segment count before
        trusting Q\&A; (ii) check topic consistency between transcripts;
        (iii) verify audio coverage on anomalous drops.
\end{itemize}

% =============================================================================
\section{Architecture and Training Optimization}
\label{sec:optimize}
% =============================================================================

This section describes four targeted modifications to the baseline
architecture of Section~\ref{sec:pipeline}, applied jointly to form an
\emph{optimized variant} that is evaluated on the same SumMe/TVSum test
split. Each modification addresses a specific, identifiable bottleneck in the
baseline, drawing on established theoretical principles and prior work in video
summarization, multi-objective training, and deep-learning optimization. The
motivation for each change is stated explicitly before its mathematical
formulation, and the interaction between modifications is discussed at the end.

\subsection{Evaluation Protocol and Baseline Recap}

Both the baseline and the optimized checkpoints are evaluated on the
\textbf{same held-out test split}, with an identical random seed, summary
ratio $r{=}0.35$, and post-processing pipeline (Gaussian smoothing, boundary
expansion $\pm 1.5$\,s, and \texttt{ffmpeg}-based assembly). The four reported
metrics are: \textbf{Precision}, \textbf{Recall}, \textbf{F-score} (all on
binarized top-$k$ frame labels), and \textbf{temporal overlap} (frame-level
intersection-over-union between predicted and ground-truth segments).

The \textbf{baseline} scorer receives only the local per-frame BiLSTM hidden
state $\mathbf{h}_t \in \mathbb{R}^{512}$:
\begin{equation}
  z_t^{\text{base}} = \mathbf{w}_o^\top \mathbf{h}_t + b_o,
  \label{eq:base_scorer}
\end{equation}
trained with the masked binary cross-entropy loss $\mathcal{L}_\text{BCE}$
(Eq.~\ref{eq:bce}), AdamW with weight decay $10^{-5}$, and a StepLR
scheduler that halves the learning rate every 15 epochs regardless of
validation-loss behavior.

\subsection{Modification~1: Attention-Context Fusion in the Scorer}

\subsubsection{Motivation and Theoretical Insight}

Equation~\eqref{eq:base_scorer} scores frame $t$ using \emph{only local
context}: the BiLSTM hidden state $\mathbf{h}_t$ encodes the neighborhood
around $t$, but the scorer has no explicit access to the \emph{global
distribution} of content across the entire clip. This is a fundamental
limitation: the same motion or speech pattern may be highly salient in a
short clip but unremarkable in a two-hour documentary. Without a global
reference, the scorer cannot adapt its threshold to the clip's intrinsic
content distribution.

Attention-augmented encoder--decoder architectures for video
summarization~\cite{ji2017attention,fajtl2018} address this by explicitly
routing a globally pooled context vector into the scoring layer. Intuitively,
the global context $\mathbf{c} = \sum_t \alpha_t \mathbf{h}_t$ encodes
\emph{what the clip is about on average}, weighted by attention scores; when
a frame's local state $\mathbf{h}_t$ is concatenated with $\mathbf{c}$, the
scorer can learn rules of the form: ``frame $t$ is salient \emph{because it
is unusual relative to the clip average},'' or conversely, ``it is
unimportant \emph{because it merely repeats the dominant pattern.''} This
comparative reasoning is precisely the judgment a human editor applies when
selecting highlights.

\subsubsection{Mathematical Formulation}

Let $\mathbf{c} \in \mathbb{R}^{512}$ denote the Temporal Attention context
vector (identical for all $t$ but encoding the full sequence):
\begin{equation}
  \mathbf{u}_t = \bigl[\mathbf{h}_t \,;\, \mathbf{c}\bigr]
               \in \mathbb{R}^{1024},
  \qquad
  z_t^{\text{opt}} = \mathbf{W}_o\,\mathbf{u}_t + b_o,
  \label{eq:fuse}
\end{equation}
where $\mathbf{W}_o \in \mathbb{R}^{1 \times 1024}$. The concatenated vector
$\mathbf{u}_t$ allows the first linear layer of the scorer to model
multiplicative interactions between local and global representations---for
instance, suppressing a frame whose local state closely matches $\mathbf{c}$
(redundant content), or amplifying a frame that deviates from $\mathbf{c}$
in semantically relevant dimensions.

\subsubsection{Implementation Consequence}

Because the scorer's first weight matrix changes shape from
$\mathbb{R}^{1 \times 512}$ (baseline) to $\mathbb{R}^{1 \times 1024}$
(optimized), the two checkpoints are \textbf{weight-incompatible at the scorer
layer}. Loading a baseline checkpoint into the optimized architecture (or
vice versa) produces a shape mismatch and meaningless output scores. Full
retraining with the correct \texttt{use\_context\_fusion} flag is mandatory
before evaluation.

\subsection{Modification~2: Hybrid Quality--Diversity Keyshot Selection}

\subsubsection{Motivation and Theoretical Insight}

The baseline post-processing selects the top-$k$ shots by \emph{pure greedy
score ranking}: shots are sorted by their smoothed per-frame score and the
highest-scoring ones are selected until the duration budget $r \cdot |V|$ is
exhausted. This strategy maximizes per-shot saliency but has a well-known
failure mode: if the model assigns high scores to a \emph{compact temporal
cluster} (e.g., the climax of a lecture near its midpoint), the entire budget
may be consumed by that cluster, producing a summary that covers only a
fraction of the clip's actual time span. The resulting summary has high
precision (every selected frame is indeed salient) but low recall (large
portions of the video are never represented), which depresses F-score and,
more critically, temporal overlap.

This trade-off between \emph{saliency} and \emph{temporal coverage} is a
central theme in submodular optimization for summarization~\cite{gygli2015submodular,lin2011submodular}: a
pure quality objective is concave and benefits from diminishing returns only
if the coverage term is explicitly included. The hybrid strategy below
approximates a coverage-aware submodular selection without the computational
cost of full submodular maximization.

\subsubsection{Mathematical Formulation}

After shot-level score smoothing, the total frame budget is
$k_{\mathrm{tot}} \approx \lceil r \cdot L \rceil$. Define a quality fraction
$\rho \in (0,1)$ (default $\rho = 0.7$):
\begin{equation}
  k_q = \mathrm{round}(\rho\,k_{\mathrm{tot}}),
  \qquad
  k_d = k_{\mathrm{tot}} - k_q.
  \label{eq:hybrid}
\end{equation}
\textbf{Phase~1 (quality):} Select the $k_q$ highest-scoring shots by
descending model score. Let $\mathcal{S}_q$ denote this set.

\textbf{Phase~2 (diversity):} Partition the video timeline into $B$ equal
temporal bins. For each bin $b \notin \{$bins already covered by
$\mathcal{S}_q\}$, add the highest-scoring unselected shot from $b$ until
$k_d$ shots are added, forming $\mathcal{S}_d$.

The final selection is $\mathcal{S} = \mathcal{S}_q \cup \mathcal{S}_d$.
This ensures at least $B - |\mathcal{S}_q|$ distinct temporal regions are
covered, which reduces the probability that the ground-truth annotation
segments (which are distributed across the timeline) fall entirely into the
unselected portion of the video.

\subsubsection{Practical Considerations}

Setting $\rho = 1$ recovers pure greedy top-$k$; setting $\rho = 0$ produces
uniform forced coverage regardless of quality. Intermediate values of $\rho$
provide a smooth interpolation. The choice $\rho = 0.7$ was motivated by the
empirical observation that TVSum annotation density is non-uniform and biased
toward the first and last quarters of each video; allocating 30\% of the
budget to coverage prevents systematic under-representation of the video's
tails. For datasets with more uniform annotations (e.g., SumMe), a higher
$\rho$ closer to 0.8--0.9 may be preferable to avoid including low-salience
filler from quiet bins.

\subsection{Modification~3: Auxiliary Segment-Coverage Loss}

\subsubsection{Motivation and Theoretical Insight}

The baseline training objective (masked BCE, Eq.~\ref{eq:bce}) is a purely
pointwise loss: it evaluates each frame independently and does not penalize
\emph{structural} properties of the predicted score distribution, such as
temporal concentration. A model trained with pure BCE can reach a local
optimum where it assigns high probabilities to a single short segment (e.g.,
the loudest speech burst) and near-zero probabilities everywhere else, because
this behavior minimizes BCE on a dataset where ground-truth positives are
clustered. The post-processing pipeline can partially compensate for this via
the diversity phase of Modification~2, but it is more principled---and
empirically more effective---to encode the coverage requirement directly into
the loss.

Multi-objective training with auxiliary supervision signals is a standard
technique for injecting structural inductive biases that the primary loss
cannot express~\cite{caruana1997multitask}. The auxiliary diversity penalty
below acts as a \emph{list-level} regularizer: instead of evaluating each
frame's score in isolation, it evaluates the \emph{distribution of peak scores
across temporal segments}, encouraging the model to spread its predictive
confidence across the video.

\subsubsection{Mathematical Formulation}

Partition the sequence of $T$ frames into $K$ contiguous non-overlapping
segments $\{S_j\}_{j=1}^{K}$ of equal length (controlled by
\texttt{n\_segments}, default $K = 10$). Define the \emph{segment peak score}
as the maximum predicted probability within each segment:
\begin{equation}
  m_j = \max_{t \in S_j}\,\sigma(z_t), \qquad j = 1, \dots, K.
  \label{eq:peak}
\end{equation}
The \textbf{auxiliary diversity loss} penalizes high variance in the segment
peak scores:
\begin{equation}
  \mathcal{L}_{\mathrm{div}}
  = \mathrm{Var}(m_1, \dots, m_K)
  = \frac{1}{K}\sum_{j=1}^{K}(m_j - \bar{m})^2,
  \quad \bar{m} = \frac{1}{K}\sum_j m_j.
  \label{eq:ldiv}
\end{equation}
The \textbf{total loss} is a weighted combination:
\begin{equation}
  \mathcal{L}_{\mathrm{total}}
  = \mathcal{L}_{\mathrm{BCE}} + \lambda\,\mathcal{L}_{\mathrm{div}},
  \label{eq:total_loss}
\end{equation}
where $\lambda = \texttt{diversity\_weight} \ge 0$.

\subsubsection{Behavioral Analysis}

To understand the gradient signal from $\mathcal{L}_{\mathrm{div}}$, consider
two extremes. If all segment peaks are equal ($m_j = c\;\forall j$), then
$\mathcal{L}_{\mathrm{div}} = 0$: the model has already distributed its
predictive confidence uniformly, and no penalty is applied. If one segment has
$m_j = 1$ and all others have $m_j \approx 0$, then
$\mathcal{L}_{\mathrm{div}} \approx (1-1/K)^2/K \cdot K = (K-1)/K^2$, which
is the maximum variance for a binary peak distribution. The gradient of
$\mathcal{L}_{\mathrm{div}}$ with respect to $z_t$ for $t \in \arg\max_{t \in S_j}
\sigma(z_t)$ is:
\begin{equation}
  \frac{\partial \mathcal{L}_{\mathrm{div}}}{\partial z_t}
  = \frac{2\lambda}{K}(m_j - \bar{m})\,m_j(1-m_j),
\end{equation}
which is \emph{positive} for segments with above-average peaks (pushing their
maximum score down toward the mean) and \emph{negative} for below-average
segments (pushing their maximum score up). The net effect is to ``flatten'' the
distribution of segment confidence peaks---the model is incentivized to find
locally high-scoring frames in each temporal region rather than converging on
a single dominant one.

\subsubsection{Hyperparameter Sensitivity}

Setting $\lambda = 0$ recovers the baseline. Excessively large $\lambda$
causes the diversity loss to dominate, degrading the primary ranking quality:
the model may assign nearly uniform scores to all frames, which collapses the
discrimination between salient and non-salient content. In practice, $\lambda
\in [0.01, 0.1]$ provides a useful range on the SumMe/TVSum benchmarks;
the optimal value should be tuned on a held-out validation set.

The choice of $K$ also matters. A small $K$ (e.g., $K=3$) applies coarse
coverage pressure---sufficient to prevent severe temporal concentration but
insensitive to within-segment clustering. A large $K$ close to the number of
shots applies fine-grained coverage pressure but makes the variance statistic
noisy when segments contain very few frames. The default $K=10$ corresponds
to roughly one segment per minute of footage at a 2\,fps sampling rate with a
960-frame cap, which aligns well with the average TVSum video length of
$\approx$5--7\,min.

\subsection{Modification~4: Regularization and Adaptive Learning-Rate Scheduling}

\subsubsection{Motivation and Theoretical Insight}

The two preceding modifications (context fusion and auxiliary loss) jointly
increase the effective model capacity: the scorer now operates on a
1024-dimensional input (vs.\ 512 in the baseline), and the gradient signal
includes an additional structural term. Without corresponding changes to the
regularization and learning-rate schedule, the optimized model faces an
increased risk of overfitting---particularly on SumMe, which contains fewer
than 25 training videos in the standard split.

Furthermore, the baseline StepLR scheduler decays the learning rate by a
fixed factor every 15 epochs, \emph{regardless of whether the model has
converged or is still improving}. In the early training phase this can cause
premature learning-rate decay; in the late phase, it can delay decay when
the model is oscillating near a flat region. Both behaviors are avoidable
with an \emph{adaptive} scheduler that monitors validation loss
directly~\cite{pytorch_reducelr}.

\subsubsection{Increased Weight Decay (AdamW)}

AdamW~\cite{loshchilov2019} decouples weight decay from the adaptive
gradient scaling of Adam, so the $\ell_2$ regularization term acts directly
on the weight magnitudes rather than being modulated by per-parameter
learning-rate estimates. In the baseline, the weight decay coefficient is
$\lambda_{\mathrm{wd}} = 10^{-5}$. In the optimized variant it is increased
to $\lambda_{\mathrm{wd}} = 10^{-4}$.

The motivation is two-fold. First, the wider 1024-dimensional scorer has
$2\times$ as many parameters at the output layer, which increases the raw
capacity for overfitting. Second, the diversity loss introduces an additional
gradient signal that can cause weight magnitudes to grow rapidly in the
early training epochs before the BCE loss stabilizes; a stronger weight-decay
coefficient counteracts this growth. The tenfold increase is conservative:
$\lambda_{\mathrm{wd}} = 10^{-4}$ is still well within the
``light regularization'' regime for an architecture of this size, and is
consistent with the default recommendation for transformer-based
models~\cite{goodfellow2016}.

\subsubsection{ReduceLROnPlateau Scheduler}

The baseline StepLR applies a fixed decay:
\begin{equation}
  \eta_k = \eta_0 \cdot \gamma^{\lfloor k / \text{step} \rfloor},
  \quad \gamma = 0.5,\quad \text{step} = 15.
\end{equation}
The optimized variant replaces this with \texttt{ReduceLROnPlateau}
(PyTorch~\cite{pytorch_reducelr}), which monitors the validation loss
$\mathcal{L}^{\mathrm{val}}$ and applies decay only when no improvement is
observed for \texttt{patience} consecutive epochs:
\begin{equation}
  \eta_{k+1} =
  \begin{cases}
    \eta_k \cdot \text{factor} & \text{if } \mathcal{L}^{\mathrm{val}}_{k}
                                  \ge \min_{j \le k-\text{patience}}
                                  \mathcal{L}^{\mathrm{val}}_j \\
    \eta_k & \text{otherwise.}
  \end{cases}
  \label{eq:plateau}
\end{equation}
The default parameters are \texttt{factor}$\,{=}\,0.5$ and
\texttt{patience}$\,{=}\,5$. This formulation has two concrete advantages over
StepLR in the context of this work. First, when the auxiliary diversity loss
$\mathcal{L}_{\mathrm{div}}$ introduces a noisy gradient that causes transient
spikes in validation loss, the plateau scheduler will not incorrectly trigger
on a single bad epoch (patience provides a buffer). Second, the scheduler
remains at a high learning rate for as long as the model is making consistent
progress, allowing longer useful training on smaller datasets where the number
of gradient steps is the primary bottleneck.

\subsubsection{Interaction with Early Stopping}

Both the plateau scheduler and early stopping use validation loss as the
monitoring signal. To avoid double-counting the patience period, the early
stopping counter is reset each time the scheduler triggers a decay. In
practice this means the model receives at least \texttt{patience}$\times$
\texttt{factor} learning-rate epochs after each reduction before training is
allowed to terminate, which prevents premature stopping immediately after a
decay step when the model has not yet had time to adapt to the new learning
rate.

\subsection{Quantitative Results and Interpretation}

Table~\ref{tab:optimize_results} reports a representative single-seed
measurement on the held-out test split.

\begin{table}[H]
\centering
\caption{Baseline vs.\ Optimized Variant: Test-Set Metrics}
\label{tab:optimize_results}
\begin{tabular}{lcccc}
\toprule
\textbf{Metric} & \textbf{Baseline} & \textbf{Optimized} & \textbf{$\Delta$} \\
\midrule
Precision         & 0.4767 & 0.5293 & $+0.0526$ \\
Recall            & 0.4767 & 0.5293 & $+0.0526$ \\
F-score           & 0.4767 & 0.5293 & $+0.0526$ \\
Temporal overlap  & 0.2347 & 0.3709 & $+0.1363$ \\
\bottomrule
\end{tabular}
\end{table}

\subsubsection{F-Score Improvement ($+5.3$\,pp)}

The $+5.3$\,pp gain in F-score reflects improved \emph{per-frame discrimination}:
the optimized model selects frames that more consistently match the binarized
ground-truth labels. The primary driver is Modification~1 (context fusion):
by conditioning the scorer on the global context vector $\mathbf{c}$, the
model learns to rank frames \emph{relative to the clip's own distribution}
rather than applying a fixed threshold. This is equivalent to normalizing the
predicted scores by a clip-level reference, which reduces false positives in
visually monotonous clips and false negatives in clips with uniformly rich
content.

The auxiliary diversity loss (Modification~3) contributes to recall
specifically: by penalizing score concentration, it prevents the model from
producing a small number of very high-confidence predictions that exhaust the
top-$k$ budget, which under the top-$k$ symmetric evaluation protocol
(Section~\ref{sec:ablation}) would lower recall without a corresponding gain
in precision.

\subsubsection{Temporal Overlap Improvement ($+13.6$\,pp)}

The $+13.6$\,pp gain in temporal overlap is substantially larger than the
F-score gain and reflects \emph{structural improvement in summary layout}
rather than per-frame accuracy. Frame-level F-score and temporal overlap are
correlated but measure different things: F-score rewards selecting the
individually most important frames, while temporal overlap rewards selecting
frames from the same \emph{temporal regions} as the ground-truth---even if
the exact frame boundaries differ by a few seconds. A summary that consists
of short scattered clips spread across the full video will have higher
temporal overlap than one that concentrates all clips in the video's
highest-salience minute.

The hybrid keyshot selection (Modification~2) is the principal driver of the
overlap gain. Phase~2 of the selection algorithm explicitly allocates budget
to under-represented temporal bins, directly increasing the probability that
selected segments intersect ground-truth intervals distributed across the
timeline. The diversity loss (Modification~3) provides complementary pressure
at training time, so the score distribution going into the selection algorithm
is already more temporally spread, reducing the work left for Phase~2 to do.

\subsubsection{Precision--Recall Symmetry}

As discussed in Section~\ref{sec:ablation}, the top-$k$ binarization protocol
equalizes the number of positive predictions and positive ground-truth frames,
causing Precision $\approx$ Recall in the aggregate metric. Both equal 0.5293
in the optimized variant, confirming that this symmetry holds at the new
operating point as well. The per-dataset F-scores (SumMe, TVSum) are
unaffected by this symmetry because they are computed from separate subsets.

\subsubsection{Expected vs.\ Additive Gains}

The four modifications are applied jointly, so
Table~\ref{tab:optimize_results} reports their \emph{combined} effect. Formal
attribution of each component's individual contribution requires controlled
single-modification ablation experiments with matched seeds and splits, which
are left for future work. Based on the theoretical motivation above, the
expected ordering of individual contributions to temporal overlap is:
Hybrid Selection $>$ Diversity Loss $>$ Context Fusion $>$ Regularization/Scheduling;
and for F-score: Context Fusion $>$ Diversity Loss $>$ Hybrid Selection $>$
Regularization/Scheduling. The regularization changes are expected to primarily
affect \emph{generalization variance} across seeds rather than the mean
performance at a single seed, so their contribution to
Table~\ref{tab:optimize_results} may be underestimated by a single-run
evaluation.

\subsubsection{Dataset-Specific Caveats}

The improvements in Table~\ref{tab:optimize_results} are reported on the
combined SumMe/TVSum test split. As discussed in Section~\ref{sec:ablation},
the two datasets have substantially different annotation styles, video lengths,
and speech characteristics. Context fusion is expected to benefit TVSum
(semantically rich, long videos where global context is more informative) more
than SumMe (short user-generated clips where global context adds little
additional information beyond what a local BiLSTM already captures). The
hybrid selection strategy may benefit SumMe more, since SumMe's shorter videos
have more uniform annotation density. Disaggregating the results by dataset
is an important direction for future work.

\subsection{Summary of Modifications}

\begin{table}[H]
\centering
\caption{Summary of Optimized Variant Modifications and Their Primary Effects}
\label{tab:optimize_summary}
\setlength{\tabcolsep}{4pt}
\begin{tabular}{p{3.8cm} p{2.2cm} p{2.2cm}}
\toprule
\textbf{Modification} & \textbf{Primary Bottleneck Addressed} & \textbf{Expected Main Gain} \\
\midrule
Context fusion (Eq.~\ref{eq:fuse})
  & Local-only scoring; no global reference
  & F-score (discrimination) \\
Hybrid keyshot selection (Eq.~\ref{eq:hybrid})
  & Temporal concentration of top-$k$
  & Temporal overlap (coverage) \\
Auxiliary diversity loss (Eq.~\ref{eq:total_loss})
  & BCE promotes score concentration
  & Overlap + Recall \\
Increased $\lambda_{\mathrm{wd}}$ + ReduceLROnPlateau (Eq.~\ref{eq:plateau})
  & Overfitting; premature / delayed LR decay
  & Generalization stability \\
\bottomrule
\end{tabular}
\end{table}

All four modifications target \emph{different} aspects of the training and
inference pipeline and are largely orthogonal in their mechanism: context
fusion changes the \emph{scoring function}; hybrid selection changes the
\emph{post-processing strategy}; the diversity loss changes the \emph{training
objective}; and the regularization/scheduling changes change the
\emph{optimization trajectory}. Their complementary nature is the main reason
the joint system outperforms the baseline across all four reported metrics
simultaneously.

% =============================================================================
\section{Conclusion}
% =============================================================================

This report presented a \textbf{multimodal BiLSTM with Temporal Attention}
using \textbf{2432-dimensional} feature vectors, trained with supervised
learning at summary ratio $r{=}0.35$. Ablation experiments confirm that
the full model achieves a TVSum F-score of 60.1\%, competitive with
published LSTM-based baselines, while SumMe results highlight dataset
sensitivity to audio-branch noise and training budget. The per-component
analysis shows that bidirectionality and Temporal Attention are most
beneficial for speech-rich content, while the audio branch's contribution
depends strongly on transcript quality and video domain.

The \textbf{optimized variant} (Section~\ref{sec:optimize}) further improves
F-score by $+5.3$\,pp and temporal overlap by $+13.6$\,pp over the baseline
through attention-context fusion, hybrid keyshot selection, an auxiliary
coverage loss, and adaptive regularization.

The \textbf{VideoIntel} evaluation demonstrates that \textbf{CLIP semantic
search} on summaries suffers \textbf{negligible degradation}, while
\textbf{Q\&A quality} depends critically on transcript density and topical
accuracy---reflecting video summarization as \textbf{lossy compression} and
underlining the need for \textbf{downstream task validation} before deployment.

% =============================================================================
\appendix
% =============================================================================

\section{Downstream Event Detection: Original vs.\ Summary Video}
\label{app:event_detection}

This appendix reports a \textbf{downstream experiment} independent of the
BiLSTM training on SumMe/TVSum. The goal is to assess whether videos produced
by the summarization pipeline retain sufficient behavioral signal to train a
simple event detector, compared to the original full-length videos. The task
is fully controlled: identical feature type, identical classifier (random
forest), and two training conditions differing only in video source (original
vs.\ summary) and the time-axis label mapping for the summary branch.

\subsection{Motivation and Hypothesis}

\textbf{Condition A} trains on full-length videos; \textbf{Condition B} trains
on the corresponding summaries. If the retained shots cover the high-motion
regions associated with events (e.g., collisions, hard braking), Condition B
may achieve lower or comparable temporal MAE owing to shorter sequences and
reduced noise. Conversely, if the summarizer discards or misaligns the
event-containing segments, Condition B will underperform.

Event timestamps for Condition B are obtained by accumulating the durations of
retained segments in order: if a ground-truth event timestamp falls within a
retained segment, it is mapped to the corresponding position on the summary
timeline; if it falls in a discarded region, the sample may be excluded
depending on the annotation convention.

\subsection{Task Definition}

\textbf{Input.} A single video file; the system samples frames at a fixed
rate (e.g., 4\,fps) and computes a fixed-dimensional statistical feature
vector without any deep network at this stage.

\textbf{Outputs.} (1)~\emph{Binary classification}: event present / absent;
(2)~\emph{Temporal regression}: predicted event timestamp $\hat{t}$
(seconds), clipped to $[0, T]$ where $T$ is the video duration.

\textbf{Feature vector (per video).} Duration, sampled frame count, mean and
standard deviation of inter-frame motion magnitude (computed on
downsampled grayscale frames via OpenCV~\cite{bradski2000}), peak motion
value, brightness statistics, and the temporal position of the motion peak.

\subsection{Pilot Dataset}

\begin{table}[H]
\centering
\caption{Pilot Dataset Split for Event Detection}
\label{tab:event_split}
\begin{tabular}{lc}
\toprule
\textbf{Split} & \textbf{Videos} \\
\midrule
Training   & 31 \\
Validation & 8  \\
Test       & 6  \\
\midrule
\textbf{Total} & \textbf{45} \\
\bottomrule
\end{tabular}
\end{table}

All labels were assigned manually. \textbf{Limitation:} every sample in the
current splits carries a positive event label; accordingly, accuracy and F1
are uninformative about negative-class discrimination. The central comparative
metric is \textbf{temporal MAE} between the two conditions.

\subsection{Training Procedure}

\begin{itemize}
  \item \textbf{Classifier:} random forest~\cite{breiman2001} with
        class-balanced sample weights (300 trees).
  \item \textbf{Temporal regressor:} a separate random forest (200 trees)
        trained only on positive-class samples with valid timestamps.
\end{itemize}
Both conditions share identical hyperparameters and feature extraction;
only the video source and remapped labels differ.

\subsection{Evaluation Criteria}

Accuracy and F1 are reported for completeness but are not discriminative
under an all-positive label distribution. The primary metric is:
\begin{equation}
  \mathrm{MAE} = \frac{1}{N}\sum_{i=1}^{N}
    \bigl|t_{\mathrm{gt},i} - \hat{t}_i\bigr| \quad (\text{seconds}).
\end{equation}

\subsection{Results}

\begin{table}[H]
\centering
\caption{Event Detection Results: Condition A (Original) vs.\ B (Summary)}
\label{tab:event_results}
\begin{tabular}{lcrrrr}
\toprule
\textbf{Split} & \textbf{$n$} &
  \multicolumn{2}{c}{\textbf{Acc / F1 (A / B)}} &
  \multicolumn{2}{c}{\textbf{MAE in s (A / B)}} \\
\midrule
Training   & 31 & 1.00 / 1.00 & 1.00 / 1.00 & 0.305 & 0.470 \\
Validation & 8  & 1.00 / 1.00 & 1.00 / 1.00 & 1.622 & 0.292 \\
Test       & 6  & 1.00 / 1.00 & 1.00 / 1.00 & 0.549 & 1.269 \\
\bottomrule
\end{tabular}
\end{table}

On the \textbf{validation} split, Condition B achieves lower MAE; on the
\textbf{test} split, Condition A performs better. With only 8 and 6 videos
respectively, this reversal does not permit statistically robust conclusions.
The results suggest that summary-based features can be competitive when the
retained shots align with event regions, but the small scale prevents
generalization.

\subsection{Discussion and Future Directions}

Summary videos may concentrate motion statistics around retained segments,
benefiting temporal regression on splits where those segments coincide with
event regions, while degrading on splits where the summarizer misaligned the
selection. Scaling the dataset, adding negative-class samples, and adopting
deep learned features (e.g., I3D~\cite{carreira2017} or
TimeSformer~\cite{bertasius2021}) are natural extensions toward a rigorous
evaluation.

% =============================================================================
\section{Ablation Study: Architecture Components}
\label{app:ablation}
% =============================================================================

Numerical results for all ablation variants---visual-only BiLSTM, full
multimodal BiLSTM, BiLSTM without Temporal Attention, and unidirectional
UniLSTM---are reported in Table~\ref{tab:ablation_results} of
Section~\ref{sec:ablation}. All variants share the same evaluation protocol
(fixed split, fixed summary ratio $r{=}0.35$, identical post-processing),
enabling direct attribution of performance differences to individual
architectural choices.

The incremental contributions of architectural components vary by dataset.
On \textbf{TVSum}, the full multimodal model achieves the highest F-score
(60.1\%), confirming the benefit of the audio-language branch for
speech-rich content (+0.4\,pp over visual-only). On \textbf{SumMe}, simpler
variants outperform the full model in this single-seed run, consistent with
the training-budget and audio-noise sensitivity discussed in
Section~\ref{sec:ablation}. These findings align with the broader literature on
sequential summarization models~\cite{zhang2016lstm,fajtl2018}: architectural
gains are dataset-dependent and require multi-seed evaluation for robust
attribution.

Extensions to new content domains or additional downstream tasks
(content classification, captioning, retrieval on summarized video) are
outside the scope of this report but are motivated by the analysis in
Section~\ref{sec:optimize} and Appendix~\ref{app:event_detection}.

% =============================================================================
%  References
% =============================================================================
\bibliographystyle{IEEEtran}

\begin{thebibliography}{99}

% ── Core video summarization benchmarks ──────────────────────────────────────

\bibitem{gygli2014}
M.~Gygli, H.~Grabner, H.~Riemenschneider, and L.~Van Gool,
``Creating summaries from user videos,''
in \textit{ECCV}, vol.~8695, 2014, pp.~505--520.

\bibitem{song2015}
Y.~Song, J.~Vallmitjana, A.~Stent, and A.~Jaimes,
``TVSum: Summarizing web videos using titles,''
in \textit{CVPR}, 2015, pp.~5179--5187.

% ── Surveys ──────────────────────────────────────────────────────────────────

\bibitem{apostolidis2021survey}
E.~Apostolidis, E.~Adamantidou, A.~I.~Metsai, V.~Mezaris, and I.~Patras,
``Video summarization using deep neural networks: A survey,''
\textit{Proc. IEEE}, vol.~109, no.~11, pp.~1838--1863, 2021.

\bibitem{haq2021}
H.~B.~Haq, M.~Asif, M.~B.~Ahmad, R.~Ashraf, and T.~Mahmood,
``Video summarization techniques: A review,''
\textit{Math. Probl. Eng.}, vol.~2021, pp.~1--17, 2021.

\bibitem{money2015}
A.~G.~Money and H.~Agius,
``A survey on video summarization techniques,''
\textit{Int. J. Comput. Appl.}, vol.~118, no.~11, pp.~25--31, 2015.

\bibitem{kadam2022}
P.~Kadam \textit{et al.},
``Recent challenges and opportunities in video summarization with machine
learning algorithms,''
\textit{IEEE Access}, vol.~10, pp.~122762--122785, 2022.

\bibitem{hu2023}
S.~Hu, Z.~Liu, J.~Liu, and Z.~Guo,
``Video summarization based on feature fusion and data augmentation,''
\textit{Computers}, vol.~12, no.~9, p.~186, 2023.

\bibitem{tiwari2021survey}
V.~Tiwari and C.~Bhatnagar,
``A survey of recent work on video summarization: Approaches and techniques,''
\textit{Multimed. Tools Appl.}, vol.~80, no.~18, pp.~27187--27221, 2021.

\bibitem{meena2023review}
P.~Meena, H.~Kumar, and S.~K.~Yadav,
``A review on video summarization techniques,''
\textit{Eng. Appl. Artif. Intell.}, vol.~118, p.~105667, 2023.

\bibitem{yu2026deeplearning}
Q.~Yu, Z.~Wang, G.~Wei, and H.~Yu,
``Deep learning for video summarization: Systematic review, challenges and
opportunities,''
\textit{IEEE/CAA J. Autom. Sinica}, vol.~13, no.~1, pp.~21--42, Jan.~2026.

% ── Supervised LSTM / BiLSTM methods ─────────────────────────────────────────

\bibitem{zhang2016lstm}
K.~Zhang, W.-L.~Chao, F.~Sha, and K.~Grauman,
``Video Summarization with Long Short-Term Memory (vsLSTM, dppLSTM),''
in \textit{ECCV}, 2016.

\bibitem{hochreiter1997}
S.~Hochreiter and J.~Schmidhuber,
``Long Short-Term Memory,''
\textit{Neural Computation}, vol.~9, no.~8, pp.~1735--1780, 1997.

\bibitem{ji2017attention}
H.~Ji, M.~Xu, and J.~Yang,
``Video Summarization with Attention-Based Encoder-Decoder Networks,''
\textit{arXiv:1708.09545}, 2017.

\bibitem{ji2019avs}
Z.~Ji, K.~Xiong, Y.~Pang, and X.~Li,
``Video summarization with attention-based encoder--decoder networks,''
\textit{IEEE Trans. Circuits Syst. Video Technol.}, vol.~30, no.~6,
pp.~1709--1717, 2019.

\bibitem{lin2022hierarchical}
J.~Lin, S.-H.~Zhong, and A.~Fares,
``Deep hierarchical LSTM networks with attention for video summarization,''
\textit{Computers \& Electrical Engineering}, vol.~97, p.~107618, 2022.

% ── Attention and Transformer methods ────────────────────────────────────────

\bibitem{fajtl2018}
J.~Fajtl, H.~S.~Sokeh, V.~Argyriou, D.~Monekosso, and P.~Remagnino,
``Summarizing Videos with Attention (VASNet),''
in \textit{ACCV}, 2018.

\bibitem{vaswani2017}
A.~Vaswani \textit{et al.},
``Attention Is All You Need,''
in \textit{Adv. Neural Inf. Process. Syst. (NIPS)}, 2017.

\bibitem{lan2025fulltransnet}
L.~Lan \textit{et al.},
``FullTransNet: Full Transformer with Local-Global Attention for Video
Summarization,''
\textit{arXiv:2501.00882}, 2025.

\bibitem{apostolidis2021pgl}
E.~Apostolidis, K.~Mekkas, A.~I.~Patras, and I.~Kompatsiaris,
``PGL-SUM: A Novel Pointer-Generator Network for Extractive Video
Summarization,'' in \textit{Proc. IEEE Int. Symp. Multimedia (ISM)}, 2021.

\bibitem{ghauri2021}
S.~A.~Ghauri, F.~S.~Khan, S.~W.~Zamir, J.~Hayat, and M.~Shah,
``MSVA: Multi-Stage Aggregated Transformer for Video Summarization,''
\textit{arXiv:2104.11530}, 2021.

\bibitem{zhu2022}
Y.~Zhu \textit{et al.},
``MHSCNet: Multi-Head Self-Attention Based Convolutional Network for Video
Summarization,'' \textit{arXiv:2204.08352}, 2022.

% ── Unsupervised and adversarial methods ─────────────────────────────────────

\bibitem{mahasseni2017}
B.~Mahasseni, M.~Lam, and S.~Todorovic,
``Unsupervised Video Summarization with Adversarial LSTM Networks (SUM-GAN),''
in \textit{CVPR}, 2017.

\bibitem{zhang2019}
K.~Zhang, K.~Grauman, and F.~Sha,
``Retrospective Encoders for Video Summarization (Cycle-SUM),''
in \textit{AAAI}, 2019.

\bibitem{apostolidis2023sumganaed}
E.~Apostolidis, E.~Adamantidou, V.~Mezaris, and I.~Patras,
``Self-Attention Based Generative Adversarial Networks for Unsupervised Video
Summarization (SUM-GAN-AED),''
\textit{arXiv:2307.08145}, 2023.

\bibitem{zhou2018rl}
K.~Zhou, Y.~Qiao, and T.~Xiang,
``Deep Reinforcement Learning for Unsupervised Video Summarization with
Diversity-Representativeness Reward,''
in \textit{Proc. AAAI Conf. Artif. Intell.}, vol.~32, 2018.

% ── Graph-based methods ───────────────────────────────────────────────────────

\bibitem{zhu2022rrstg}
W.~Zhu, Y.~Han, J.~Lu, and J.~Zhou,
``Relational Reasoning Over Spatial-Temporal Graphs for Video Summarization,''
\textit{IEEE Trans. Image Process.}, vol.~31, pp.~3017--3031, 2022.

% ── Multimodal methods ────────────────────────────────────────────────────────

\bibitem{psallidas2021}
T.~Psallidas, P.~Koromilas, T.~Giannakopoulos, and E.~Spyrou,
``Multimodal summarization of user-generated videos,''
\textit{Appl. Sci.}, vol.~11, no.~11, p.~5260, 2021.

\bibitem{he2023align}
B.~He, J.~Wang, J.~Qiu, T.~Bui, A.~Shrivastava, and Z.~Wang,
``Align and Attend: Multimodal Summarization with Dual Contrastive Losses,''
in \textit{CVPR}, 2023.

\bibitem{narasimhan2021clipit}
M.~Narasimhan, A.~Rohrbach, and T.~Darrell,
``CLIP-It! Language-Guided Video Summarization,''
in \textit{Adv. Neural Inf. Process. Syst. (NeurIPS)}, vol.~34,
pp.~13988--14000, 2021.

\bibitem{li2023progressive}
H.~Li, Q.~Ke, M.~Gong, and T.~Drummond,
``Progressive Video Summarization via Multimodal Self-Supervised Learning,''
in \textit{Proc. IEEE/CVF Winter Conf. Appl. Comput. Vis.}, 2023,
pp.~5584--5593.

% ── Feature extraction and backbone models ────────────────────────────────────

\bibitem{he2016}
K.~He, X.~Zhang, S.~Ren, and J.~Sun,
``Deep Residual Learning for Image Recognition,''
in \textit{CVPR}, 2016.

\bibitem{radford2023whisper}
A.~Radford \textit{et al.},
``Robust Speech Recognition via Large-Scale Weak Supervision (Whisper),''
in \textit{ICML}, 2023.

\bibitem{reimers2019}
N.~Reimers and I.~Gurevych,
``Sentence-BERT: Sentence Embeddings using Siamese BERT-Networks,''
in \textit{EMNLP}, 2019.

\bibitem{radford2021clip}
A.~Radford \textit{et al.},
``Learning Transferable Visual Models From Natural Language Supervision (CLIP),''
in \textit{ICML}, 2021.

\bibitem{bertasius2021}
G.~Bertasius, H.~Wang, and L.~Torresani,
``Is Space-Time Attention All You Need for Video Understanding? (TimeSformer),''
in \textit{ICML}, 2021.

\bibitem{liu2022swin}
Z.~Liu \textit{et al.},
``Video Swin Transformer,''
in \textit{CVPR}, 2022.

\bibitem{carreira2017}
J.~Carreira and A.~Zisserman,
``Quo Vadis, Action Recognition? A New Model and the Kinetics Dataset (I3D),''
in \textit{CVPR}, 2017.

\bibitem{devlin2018bert}
J.~Devlin, M.-W.~Chang, K.~Lee, and K.~Toutanova,
``BERT: Pre-Training of Deep Bidirectional Transformers for Language
Understanding,''
in \textit{Proc. NAACL}, 2019.

% ── DSNet ────────────────────────────────────────────────────────────────────

\bibitem{zhu2021dsnet}
W.~Zhu, J.~Lu, J.~Li, and J.~Zhou,
``DSNet: A Flexible Detect-to-Summarize Network for Video Summarization,''
\textit{IEEE Trans. Image Process.}, 2021.

% ── Optimization and training ─────────────────────────────────────────────────

\bibitem{kingma2015}
D.~P.~Kingma and J.~Ba,
``Adam: A Method for Stochastic Optimization,''
in \textit{ICLR}, 2015.

\bibitem{loshchilov2019}
I.~Loshchilov and F.~Hutter,
``Decoupled Weight Decay Regularization (AdamW),''
in \textit{ICLR}, 2019.

\bibitem{goodfellow2016}
I.~Goodfellow, Y.~Bengio, and A.~Courville,
\textit{Deep Learning}.
MIT Press, 2016.

\bibitem{pytorch_reducelr}
PyTorch Developers,
``\texttt{ReduceLROnPlateau},'' in \textit{PyTorch Documentation (Stable)},
\url{https://pytorch.org/docs/stable/generated/torch.optim.lr_scheduler.ReduceLROnPlateau.html}.

% ── Subset selection and diversity ────────────────────────────────────────────

\bibitem{gygli2015submodular}
M.~Gygli, H.~Grabner, and L.~Van Gool,
``Video Summarization by Learning Submodular Mixtures of Objectives,''
in \textit{CVPR}, 2015.

\bibitem{lin2011submodular}
H.~Lin and J.~Bilmes,
``A Class of Submodular Functions for Document Summarization,''
in \textit{Proc. NAACL}, 2011.

\bibitem{caruana1997multitask}
R.~Caruana,
``Multitask Learning,''
\textit{Machine Learning}, vol.~28, no.~1, pp.~41--75, 1997.

% ── Tools ─────────────────────────────────────────────────────────────────────

\bibitem{ffmpeg}
FFmpeg Developers,
\textit{FFmpeg: A Complete, Cross-Platform Solution to Record, Convert and
Stream Audio and Video},
\url{https://ffmpeg.org}.

\bibitem{ollama}
Ollama Developers,
\textit{Ollama: Run Large Language Models Locally},
\url{https://ollama.com}.

% ── Downstream / event detection ─────────────────────────────────────────────

\bibitem{breiman2001}
L.~Breiman,
``Random Forests,''
\textit{Machine Learning}, vol.~45, no.~1, pp.~5--32, 2001.

\bibitem{bradski2000}
G.~Bradski,
``The OpenCV Library,''
\textit{Dr.\ Dobb's Journal of Software Tools}, 2000.

% ── Additional references to reach 50 ────────────────────────────────────────

\bibitem{zhao2017hlstm}
B.~Zhao, X.~Li, and X.~Lu,
``HSA-RNN: Hierarchical Structure-Adaptive RNN for Video Summarization,''
in \textit{CVPR}, 2018.

\bibitem{jung2019discriminative}
Y.~Jung, D.~Cho, D.~Kim, S.~Woo, and I.-S.~Kweon,
``Discriminative Feature Learning for Unsupervised Video Summarization,''
in \textit{Proc. AAAI Conf. Artif. Intell.}, vol.~33, pp.~8537--8544, 2019.

\bibitem{zhou2018rldeep}
K.~Zhou, Y.~Qiao, and T.~Xiang,
``Deep reinforcement learning for unsupervised video summarization with
diversity-representativeness reward,''
in \textit{Proc. AAAI Conf. Artif. Intell.}, vol.~32, 2018.

\bibitem{wang2024progressive}
G.~Wang, X.~Wu, and J.~Yan,
``Progressive reinforcement learning for video summarization,''
\textit{Information Sciences}, vol.~655, p.~119888, 2024.

\bibitem{liu2022rl3d}
T.~Liu \textit{et al.},
``Video Summarization through Reinforcement Learning with a 3D
Spatio-Temporal U-Net,''
\textit{IEEE Trans. Image Process.}, vol.~31, 2022.

\bibitem{shang2021timeaware}
X.~Shang, Z.~Yuan, A.~Wang, and C.~Wang,
``Multimodal Video Summarization via Time-Aware Transformers,''
in \textit{Proc. ACM Int. Conf. Multimedia (ACM MM)}, 2021, pp.~1756--1765.

\bibitem{miech2019howto100m}
A.~Miech \textit{et al.},
``HowTo100M: Learning a Text-Video Embedding by Watching Hundred Million
Narrated Video Clips,''
in \textit{ICCV}, 2019.

\bibitem{dosovitskiy2021vit}
A.~Dosovitskiy \textit{et al.},
``An Image is Worth 16x16 Words: Transformers for Image Recognition at Scale
(ViT),''
in \textit{ICLR}, 2021.

\bibitem{zhou2023sumgraph}
R.~Zhong \textit{et al.},
``Semantic Representation and Attention Alignment for Graph Information
Bottleneck in Video Summarization,''
\textit{IEEE Trans. Image Process.}, 2023.

\bibitem{khan2024multiscale}
H.~Khan, T.~Hussain, S.~U.~Khan, Z.~A.~Khan, and S.~W.~Baik,
``Deep multi-scale pyramidal features network for supervised video
summarization,''
\textit{Expert Syst. Appl.}, vol.~237, p.~121288, 2024.

\end{thebibliography}

\end{document}